% R590: EGA I, Chapter0, whole3.2.6 including printed28 continuation.
% Traducción: OpenAI Codex, GPT-6.1 Sol, esfuerzo Ultra.
\begin{env}[3.2.6]
\label{0.3.2.6-fr}
Supongamos siempre que $\boldsymbol{K}$ admite límites proyectivos. Entonces la
categoría de los \emph{haces sobre $X$ con valores en $\boldsymbol{K}$} también admite
\emph{límites proyectivos}; si $(\mathcal{F}_\lambda)$ es un sistema
proyectivo de haces sobre $X$ con valores en $\boldsymbol{K}$, las fórmulas
\[
  \mathcal{F}(U)=\varprojlim_\lambda\mathcal{F}_\lambda(U)
\]
definen, en efecto, un prehaz con valores en $\boldsymbol{K}$, y la
verificación del axioma (F) se deduce de la transitividad de los límites
proyectivos; es inmediato que $\mathcal{F}$ es entonces el límite proyectivo
de los $\mathcal{F}_\lambda$.

Cuando $\boldsymbol{K}$ es la categoría de los conjuntos, para todo sistema
proyectivo $(\mathfrak{H}_\lambda)$ tal

\oldpage[0\textsubscript{I}]{28}

que $\mathfrak{H}_\lambda$ sea un \emph{subhaz} de
$\mathcal{F}_\lambda$ para todo $\lambda$,
$\varprojlim\mathfrak{H}_\lambda$ se identifica canónicamente con un
\emph{subhaz} de $\varprojlim\mathcal{F}_\lambda$. Si
$\boldsymbol{K}$ es la categoría de los grupos conmutativos, el funtor covariante
$\varprojlim\mathcal{F}_\lambda$ es \emph{aditivo} y \emph{exacto por la izquierda}.
\end{env}
